\documentclass[10pt,a4paper]{amsart}
\usepackage[margin=19mm]{geometry}
\usepackage{amsmath,amssymb,mathtools}
\usepackage[scr=rsfs,cal=euler]{mathalfa}
\usepackage{xfrac}
\usepackage[unicode,hidelinks,bookmarks=false]{hyperref}
\newcommand{\ie}{i.e.\ }
\newcommand{\cf}{cf.\ }
\newcommand{\resp}{respectively\ }
\newcommand{\Subsection}{\subsection}
\providecommand{\nobreakdash}{\nobreak\hbox{-}}
\setlength{\emergencystretch}{2em}
\multlinegap0pt
\usepackage{amssymb}
\usepackage{tikz}
\usepackage{xcolor}
\newtheorem{lemma}{Lemma}
\newtheorem{theorem}{Theorem}
\title{\textbf{Extremal Mostar Index of Graphs with Given Number of Cut Edges}}
\author{Sunilkumar M. Hosamani}
\date{Source: arXiv:2604.06839v1 (CC BY 4.0)}
\begin{document}
\maketitle

\noindent\emph{Source and licence.} \textbf{Extremal Mostar Index of Graphs with Given Number of Cut Edges}. Version arXiv:2604.06839v1 (CC BY 4.0), distributed by arXiv under the Creative Commons Attribution 4.0 International licence, http://creativecommons.org/licenses/by/4.0/, as recorded in the arXiv licence metadata for this exact version and shown on the abstract page https://arxiv.org/abs/2604.06839v1. Full licence notice: \texttt{licenses/arxiv-2604.06839-CC-BY-4.0.txt}.\par\smallskip

\noindent\textit{Editorial scope.} Counted unit: the article's theorem thm:max (source0.tex lines 157--163). The article's complete original proof of it is delivered with it (source0.tex lines 165--185, one \texttt{proof} environment of 21 lines). No mathematics is written, repaired or re-derived: the statement and the proof are the article's own lines, in the article's own notation, and no retained line is substituted or re-flowed.  Retained uncounted as the source-local context the proof needs: lines 73--78; lines 95--100; lines 105--107; lines 119--121.  Nothing was omitted from any retained span, so the package carries no omission record.  No retained line contains a citation, so the wrapper carries no bibliography and no bibliography entry.  No retained line contains a figure environment, an image-inclusion command or drawing code, so no image is delivered and the package declares no asset dependency.  Editorial formatting only, in addition: an AMS \texttt{amsart} preamble that restates the article's own declarations and macros used by the retained lines, namely the environments \texttt{lemma}, \texttt{theorem} on separate counters, together with the standard packages those lines need.  The retained environments print in this document as Lemma 1, Theorem 1 in order of appearance, whereas the article prints the same items with the numbering of its own full document.  The complete native source of the article as distributed in the arXiv e-print for this exact version is bundled unchanged as \texttt{sources/198046/source0.tex}.
\begin{lemma}\label{lem:pendant}
If \(uv\) is a pendant edge with \(d_G(u)=1\) and \(d_G(v)\ge 2\), then
\[
|n_u(uv)-n_v(uv)| = n-2.
\]
\end{lemma}

\begin{lemma}\label{lem:contract}
Let \(G\) be a connected graph and let \(uv\) be a non‑pendant cut edge. Construct \(G'\) by contracting \(uv\) into a single vertex \(u\) (deleting \(v\) and identifying it with \(u\)), then adding a new pendant vertex \(v'\) adjacent only to \(u\). Then \(G'\) has the same number of cut edges as \(G\), and
\[
Mo(G') > Mo(G).
\]
\end{lemma}

\begin{lemma}\label{lem:movependant}
Let \(G\) be a connected graph with at least two non‑pendant vertices \(x,y\) such that \(d_G(x)\ge d_G(y)\ge 2\). Suppose \(p\) is a pendant neighbour of \(y\). Form \(G'\) by deleting edge \(py\) and adding edge \(px\) (i.e., move the leaf \(p\) from \(y\) to \(x\)). Then \(Mo(G') > Mo(G)\).
\end{lemma}

\begin{lemma}\label{lem:addedge}
Let \(G\) be a connected graph and let \(u,v\) be non‑adjacent vertices. Then \(Mo(G+uv) > Mo(G)\).
\end{lemma}

\begin{theorem}\label{thm:max}
Let \(G\) be a connected graph of order \(n\ge 2\) with exactly \(k\) cut edges, where \(1\le k\le n-1\). Then
\[
Mo(G) \le k(n-2) + (n-k-1)k.
\]
Equality holds if and only if \(G\cong K_{n-k}^k\).
\end{theorem}

\begin{proof}
Let \(G\) be a graph achieving the maximum Mostar index among all connected graphs of order \(n\) with exactly \(k\) cut edges.

\textbf{Step 1. All cut edges are pendant.} If there were a non‑pendant cut edge, Lemma~\ref{lem:contract} would produce another graph with the same number of cut edges and a strictly larger Mostar index, contradicting maximality. Hence every cut edge of \(G\) is incident to a leaf.

\textbf{Step 2. All pendant edges are attached to a single vertex.} Let \(L\) be the set of leaves (pendant vertices). Each leaf is incident to exactly one cut edge. Let \(X\) be the set of internal vertices incident to these cut edges. If \(|X|\ge2\), pick two vertices \(x,y\in X\) with \(d(x)\ge d(y)\). By Lemma~\ref{lem:movependant}, moving a leaf from \(y\) to \(x\) increases the Mostar index while preserving the number of cut edges (the moved edge remains a cut edge). Repeating this process, we can transfer all leaves to the vertex with the largest degree without decreasing the index. Thus in a maximal graph, all pendant edges share a common endpoint, say \(v_0\).

\textbf{Step 3. The subgraph induced by the non‑pendant vertices is complete.} Let \(H = G - L\) be the subgraph obtained by deleting all leaves. Then \(H\) has \(n-k\) vertices. If two vertices \(u,w\in V(H)\) are non‑adjacent, adding the edge \(uw\) (Lemma~\ref{lem:addedge}) increases \(Mo(G)\) and does not create new cut edges (since edges inside \(H\) are not bridges). This would contradict maximality. Hence \(H\) is a complete graph \(K_{n-k}\).

\textbf{Step 4. Identify the extremal graph.} The structure is now forced: \(G\) consists of a complete graph on \(n-k\) vertices, with one distinguished vertex \(v_0\) that is additionally adjacent to \(k\) pendant leaves. This graph is exactly \(K_{n-k}^k\).

\textbf{Step 5. Compute the Mostar index of \(K_{n-k}^k\).} Label the hub as \(v_0\) and the other clique vertices as \(w_1,\dots,w_{n-k-1}\). The leaves are \(p_1,\dots,p_k\) attached to \(v_0\). 
\begin{itemize}
\item For each pendant edge \(v_0p_i\): \(b = n-2\) (Lemma~\ref{lem:pendant}). Total from pendant edges: \(k(n-2)\).
\item For each edge \(v_0w_j\): The side of \(v_0\) contains \(v_0\) itself and all \(k\) leaves, so \(n_{v_0}=1+k\). The side of \(w_j\) contains only \(w_j\) (all other vertices are equidistant or closer to \(v_0\)? Check: any other \(w_\ell\) is at distance 1 from both \(v_0\) and \(w_j\); leaves are distance 2 from \(w_j\) but 1 from \(v_0\), so they are strictly closer to \(v_0\). Hence \(n_{w_j}=1\)). Thus \(b = (1+k)-1 = k\). There are \(n-k-1\) such edges, contributing \(k(n-k-1)\).
\item For any edge \(w_j w_\ell\) with \(j\neq\ell\): The graph is symmetric under swapping \(w_j\) and \(w_\ell\) (fixing \(v_0\) and all leaves), so \(n_{w_j}=n_{w_\ell}\) and \(b=0\). There are \(\binom{n-k-1}{2}\) such edges, contributing 0.
\end{itemize}
Summing gives \(Mo(K_{n-k}^k)=k(n-2)+k(n-k-1)\).

Thus the upper bound holds, and equality forces \(G\cong K_{n-k}^k\). This completes the proof of Theorem~\ref{thm:max}.
\end{proof}

\end{document}
